\chapter*{Introduction générale}
\addcontentsline{toc}{chapter}{Introduction générale}
\argossetchapter{Introduction générale}

Aujourd’hui, le transport maritime constitue la colonne vertébrale de l’économie globale, assurant plus de 80\,\% du volume du commerce mondial, d'après le rapport \textit{Review of Maritime Transport 2025} de la CNUCED\cite{unctad2025}. En 2025, le volume des échanges conteneurisés a atteint un niveau historique de près de 192,9~millions d’EVP, selon Container Trades Statistics (CTS)\cite{cts2025}. Cette dynamique repose sur un outil standardisé~: le conteneur maritime, clé de voûte de la chaîne logistique multimodale, dont le marché mondial devrait croître de 4,1\,\% par an jusqu’en 2035, selon Global Market Insights (2025)\cite{gmi2025}.

Pourtant, derrière cette robustesse mécanique se cache une vulnérabilité majeure~: une fois scellé et chargé, le conteneur devient un véritable «~géant aveugle~». Si les systèmes d’information portuaires permettent une traçabilité géographique globale, le suivi en temps réel de l’intégrité physique de la marchandise reste un défi complexe, la majorité des conteneurs (\textit{dry}) étant des unités passives, sans alimentation ni connectivité embarquées, contrairement aux \textit{reefers}, equiper chaque unité d’un capteur autonome pose des contraintes de coût, d’autonomie énergétique et de couverture réseau en haute mer. Durant des voyages de plusieurs semaines, les parties prenantes manquent ainsi d’une visibilité cruciale. En cas d’avarie, de variation thermique ou de rupture de scellé, le constat du sinistre s'effectue fréquemment à la destination finale, entraînant une dilution des responsabilités. La détermination de l'instant précis de l'incident devient alors une épreuve probatoire complexe, justifiant toute l'importance des mécanismes de présomption de responsabilité du transporteur mis en lumière par Frederick Babu et Aimé Mbuangi-Mavungu dans le cadre normatif de l'OHADA et de la CEMAC\cite{babu_mbuangi2023}.

Cette opacité engendre un flou réglementaire que les nouvelles réglementations internationales viennent durcir. Depuis le 1\textsuperscript{er}~janvier 2026, les amendements aux conventions SOLAS (règle V/31-32) et MARPOL (Protocole I art. V) de l’OMI\cite{omi2026} obligent les navires à déclarer sans délai toute perte de conteneur en mer, rendant une traçabilité active indispensable.

Cet enjeu prend une dimension critique en Afrique centrale, et particulièrement au Cameroun. Le Port Autonome de Douala concentre un marché conteneurisé national de 183\,956~EVP en 2025\cite{pad2025}, tout en assurant le transit vital vers le Tchad et la République Centrafricaine, dont respectivement environ 80\,\% et 90\,\% des importations dépendent des ports camerounais\cite{corridor2025}. Le complexe portuaire de Kribi s’affirme quant à lui avec une croissance fulgurante, atteignant 555\,398~EVP en 2025, soit un bond de 82\,\% en un an porté par la phase II de son terminal\cite{pak2026}. Or, malgré ces volumes et la modernisation des infrastructures, l’intégrité de la cargaison reste hors du contrôle direct des opérateurs logistiques, le caractère scellé des conteneurs empêchant toute inspection sans rompre l’intégrité douanière.

\textbf{D’où la question de savoir comment garantir la traçabilité continue et l’intégrité physique vérifiable d’une cargaison enfermée dans un environnement métallique hostile et fréquemment isolé des réseaux, tout en produisant une preuve fiable et conforme aux exigences de l’OMI en cas d’avarie ou de litige ?}

Pour transformer ce maillon aveugle en nœud d'information transparent, ce projet vise à développer une solution embarquée de nouvelle génération : une boîte noire logistique, inspirée des enregistreurs de vol aéronautiques et adaptée aux contraintes du transport en Afrique centrale. Ce dispositif devra collecter, horodater et transmettre de manière autonome des données critiques (température, humidité, chocs, ouvertures, géolocalisation)

Le présent document s'articule autour de trois parties. La première, « Cadre opératoire du projet », dresse un état des lieux du transport conteneurisé et identifie les failles actuelles en matière de traçabilité, à travers la présentation du projet et l'étude de marché. La deuxième, « Évaluation du projet », se concentre sur les études de faisabilité et de viabilité, et démontre la pertinence de l'initiative sur les plans technique, juridique, opérationnel et financier. La troisième, « Étude technique et implémentation de la solution », détaille l'ingénierie de la boîte noire ARGOS~: architecture matérielle, protocoles de communication et mise en œuvre concrète du dispositif.
